\documentclass[10pt,a4paper]{amsart}
\usepackage{amsmath,amssymb,amsthm,graphicx}
\usepackage{verbatim}
\newenvironment{reference}{\comment}{\endcomment}
%\newenvironment{reference}{}{}
\newenvironment{slogan}{\comment}{\endcomment}
\newenvironment{history}{\comment}{\endcomment}

% For commutative diagrams we use Xy-pic
\usepackage[all]{xy}

% We use 2cell for 2-commutative diagrams.
\xyoption{2cell}
% Theorem environments.
%
\theoremstyle{plain}
\newtheorem{theorem}[subsection]{Theorem}
\newtheorem{proposition}[subsection]{Proposition}
\newtheorem{lemma}[subsection]{Lemma}

\theoremstyle{definition}
\newtheorem{definition}[subsection]{Definition}
\newtheorem{example}[subsection]{Example}
\newtheorem{exercise}[subsection]{Exercise}
\newtheorem{situation}[subsection]{Situation}

\theoremstyle{remark}
\newtheorem{remark}[subsection]{Remark}
\newtheorem{remarks}[subsection]{Remarks}

\numberwithin{equation}{subsection}

% Macros
%
\def\lim{\mathop{\mathrm{lim}}\nolimits}
\def\colim{\mathop{\mathrm{colim}}\nolimits}
\def\Spec{\mathop{\mathrm{Spec}}}
\def\Hom{\mathop{\mathrm{Hom}}\nolimits}
\def\Ext{\mathop{\mathrm{Ext}}\nolimits}
\def\SheafHom{\mathop{\mathcal{H}\!\mathit{om}}\nolimits}
\def\SheafExt{\mathop{\mathcal{E}\!\mathit{xt}}\nolimits}
\def\Sch{\mathit{Sch}}
\def\Mor{\mathop{\mathrm{Mor}}\nolimits}
\def\Ob{\mathop{\mathrm{Ob}}\nolimits}
\def\Sh{\mathop{\mathit{Sh}}\nolimits}
\def\NL{\mathop{N\!L}\nolimits}
\def\CH{\mathop{\mathrm{CH}}\nolimits}
\def\proetale{{pro\text{-}\acute{e}tale}}
\def\etale{{\acute{e}tale}}
\def\QCoh{\mathit{QCoh}}
\def\Ker{\mathop{\mathrm{Ker}}}
\def\Im{\mathop{\mathrm{Im}}}
\def\Coker{\mathop{\mathrm{Coker}}}
\def\Coim{\mathop{\mathrm{Coim}}}

% Boxtimes
%
\DeclareMathSymbol{\boxtimes}{\mathbin}{AMSa}{"02}

%
% Macros for moduli stacks/spaces
%
\def\QCohstack{\mathcal{QC}\!\mathit{oh}}
\def\Cohstack{\mathcal{C}\!\mathit{oh}}
\def\Spacesstack{\mathcal{S}\!\mathit{paces}}
\def\Quotfunctor{\mathrm{Quot}}
\def\Hilbfunctor{\mathrm{Hilb}}
\def\Curvesstack{\mathcal{C}\!\mathit{urves}}
\def\Polarizedstack{\mathcal{P}\!\mathit{olarized}}
\def\Complexesstack{\mathcal{C}\!\mathit{omplexes}}
% \Pic is the operator that assigns to X its picard group, usage \Pic(X)
% \Picardstack_{X/B} denotes the Picard stack of X over B
% \Picardfunctor_{X/B} denotes the Picard functor of X over B
\def\Pic{\mathop{\mathrm{Pic}}\nolimits}
\def\Picardstack{\mathcal{P}\!\mathit{ic}}
\def\Picardfunctor{\mathrm{Pic}}
\def\Deformationcategory{\mathcal{D}\!\mathit{ef}}
\usepackage{cleveref}
\crefname{originalreference}{original reference}{original references}
\Crefname{originalreference}{Original reference}{Original references}
\makeatletter
\newlabel{lemma-localization}{{096V}{source}{}{sourcepointer-1}{}}
\newlabel{lemma-localization@cref}{{[originalreference][0][]096V}{[1][0][]source}{}{}{}}
\makeatother
\title{Retractions of faithfully flat etale maps persist after localizing along a closed subscheme}
\author{The Stacks Project Authors}
\date{stacks-project:proetale.tex:local-snapshot-20261012; GNU FDL 1.2 or later}
\begin{document}
\maketitle
\noindent\textit{Editorial note.} All delivered mathematical text is the author's own, copied from the pinned Stacks Project chapter file proetale.tex; no statement or proof step was rewritten, reordered or removed.\par\medskip
\noindent\textit{Editorial note.} The counted proof is the author's complete argument: the descent of the given etale algebra over the localization to an etale algebra over A, the adjunction of a localization making it faithfully flat so that the hypothesis applies, and the localization of the resulting retraction back to A\_Z\textasciicircum{}tilde.\par\medskip
\noindent\textit{Editorial note.} The label lemma-have-sections-localize names the counted statement, and the proof environment immediately following it in the source is delivered verbatim as its proof.\par\medskip
\noindent\textit{Editorial note.} The counted statement is carried by the author's own lemma environment in the chapter source.\par\medskip
\noindent\textit{Editorial note.} Pre-open gate: python operations/agent-authoring/tools/precard\_gate.py --source proetale.tex --label lemma-have-sections-localize --start 1702 --end 1720 --local-id 28008 returned RESULT=OK (L28008 free, occupancy FREE, public interval clear). Candidate taken from sources/stacks-proetale/candidates-halfB.json.\par\medskip
\noindent\textit{Editorial note.} The statement and proof are copied from the pinned chapter file; the cited inputs are kept as external pointers with their canonical Stacks tags.\par\medskip
\noindent\textit{Allowed standard prerequisites (editorial).}\begin{itemize}
\item Descent of an etale algebra along the localization map. Algebra, the lemma describing etale algebras over a localization as the base change of etale algebras over the original ring. Supplies the etale ring map A -> B' with C = B' tensor over A with A\_Z\textasciicircum{}tilde.
\item Openness of the image of an etale morphism and covers by the image and a basic open. Algebra, the proposition that the image of an etale (more generally flat finitely presented) morphism is open, and the elementary cover by a basic open D(f). Supply the element f in the ideal such that the spectrum of A is the union of the image U' and D(f).
\item The localization construction attached to a closed subscheme. Same chapter, the construction of A\_Z\textasciicircum{}tilde together with the fact that Z is a closed subscheme of its spectrum. Supplies the localization along which the obtained retraction is taken.
\end{itemize}\medskip
\ifdefined\hypertarget\hypertarget{sourcepointer-1}{}\fi
\noindent\textit{Source reference.} Original reference 096V: sources/stacks-proetale/pinned-original-source.tex, source lines 710--710. Prerequisite pointer invoked at proof line(s) 1706; it names a standard fact of the same chapter that the counted argument uses, and every step of the counted proof built on it is retained verbatim in the delivered body.\par\medskip
\begin{lemma}
\label{lemma-have-sections-localize}
Let $A$ be a ring such that every faithfully flat \'etale ring map
$A \to B$ has a retraction. Let $Z \subset \Spec(A)$ be a closed subscheme.
Let $A \to A_Z^\sim$ be as constructed in Lemma \ref{lemma-localization}.
Then every faithfully flat \'etale ring map $A_Z^\sim \to C$ has
a retraction.
\end{lemma}

\begin{proof}
There exists an \'etale ring map $A \to B'$ such that
$C = B' \otimes_A A_Z^\sim$ as $A_Z^\sim$-algebras.
The image $U' \subset \Spec(A)$ of $\Spec(B') \to \Spec(A)$
is open and contains $V(I)$, hence we can find $f \in I$ such
that $\Spec(A) = U' \cup D(f)$. Then $A \to B' \times A_f$
is \'etale and faithfully flat. By assumption there is a retraction
$B' \times A_f \to A$. Localizing we obtain the desired retraction
$C \to A_Z^\sim$.
\end{proof}

\end{document}
